\loai{Số neutron được tạo ra}
\begin{vidu}%[3P7-5.2B][Loai1]
Sự phân hạch của hạt nhân uranium $\Hn{92}{235}{U}$ khi hấp thụ một neutron chậm xảy ra theo nhiều cách. Một trong các cách đó được cho bởi phương trình $\Hn{92}{235}{U}+\Hn{0}{1}{n}\to  \Hn{54}{140}{Xe}+\Hn{38}{94}{Sr}+k\Hn{0}{1}{n}$. Số neutron được tạo ra trong phản ứng này là
\choice
{$k = 3$}
{$k = 6$}
{$k = 4$}
{\True $k = 2$}
\loigiai{
\begin{itemize}
\item $\Hn{92}{235}{U}+\Hn{0}{1}{n}\to  \Hn{54}{140}{Xe}+\Hn{38}{94}{Sr}+k\Hn{0}{1}{n}$.
\item Áp dụng định luật bảo toàn số nuclôn ta có: $235+1=140+94+k\Rightarrow k=2$.
\end{itemize}
}
\end{vidu}
\loai{Tính năng lượng của 1 phân hạch}
\begin{vidu}%[3P7-5.2B][Loai1]
Xét phản ứng phân hạch uranium $^{235}U$ có phương trình: $\Hn{92}{235}{U}+n\Rightarrow \Hn{42}{95}{Mo}+\Hn{57}{139}{La}+2n+7e^-$. Cho biết
${m_U}=234,99u;\ m_{Mo}=94,88u;\ m_{La}=138,87u;\ {m_n}=1,0087u$. Biết $1uc^2=931,5\ MeV$ và bỏ qua khối lượng electron. Năng lượng mà một phân hạch toả ra là
\choice
{107 MeV}
{\True 215,5 MeV}
{234 MeV}
{206 MeV}
\loigiai{
Năng lượng tỏa ra trong phản ứng là:
\begin{align*}
W=\left( {m_U}+{m_n} \right)-\left( m_{Mo}+m_{La}+2{m_n} \right). c^2=215,5\ MeV.
\end{align*}
}
\end{vidu}
\loai{Tính năng lượng khi phân hạch hết m gam}
\begin{vidu}%[3P7-5.2B][Loai1]
\nguon{QG - 2017} Cho rằng một hạt nhân Uranium \Hn{92}{235}{U} khi phân hạch thì tỏa ra năng lượng là 200 MeV. Lấy $N_A=6,{02.10}^{23}\ mo{l}^{- 1}$, $1\ eV=1,{6.10}^{- 19}\ J$ và khối lượng mol của Uranium \Hn{92}{235}{U} là 235 $g/mol$. Năng lượng tỏa ra khi 2 g Uranium \Hn{92}{235}{U} phân hạch hết là
\choice
{$9,{6.10}^{10}$ J}
{$10,{3.10}^{23}$ J}
{$16,{4.10}^{23}$ J}
{\True $16,{4.10}^{10}$ J}
\loigiai{
Năng lượng tỏa ra khi phân hạch hết 2 g Urani: $Q=N.W=\dfrac{m}{A}N_A.W=1,64.10^{11}\ J$.
}
\end{vidu}
\loai{Tính khối lượng nhiên liệu}
\begin{vidu}%[3P7-5.2G][Loai1]
\nguon{QG - 2017}  Giả sử một nhà máy điện hạt nhân dùng nhiên liệu Uranium \Hn{92}{235}{U}. Biết công suất phát điện là 500 MW và hiệu suất chuyển hóa năng lượng hạt nhân thành điện năng là $20\%$. Cho rằng khi một hạt nhân Uranium \Hn{92}{235}{U} phân hạch thì tỏa ra năng lượng là $3,{2.10}^{- 11}$ J. Lấy $N_A=6,{02.10}^{23}\ mol^{- 1}$ và khối lượng mol của \Hn{92}{235}{U} là $235\ g/mol$. Nếu nhà máy hoạt động liên tục thì lượng Uranium \Hn{92}{235}{U} mà nhà máy cần dùng trong 365 ngày là
\choice
{1352,5 kg}
{\True 962 kg}
{1121 kg}
{1421 kg}	
\loigiai{
\begin{itemize}
\item Năng lượng hạt nhân cần cho 365 ngày: $Q=\dfrac{P_c.t}{H}$.
\item Năng lượng của m gam hạt nhân phân hạch hết: $Q'=NW=\dfrac{m}{M}.N_A.W$.
\item Lượng Uranium cần dùng: $Q=Q' \Leftrightarrow \dfrac{P_c.t}{H}=\dfrac{mN_AW}{M} \Rightarrow m=\dfrac{MP_ct}{HN_AW}=961763\ g\approx 962\ kg$.
\end{itemize}
}
\end{vidu}
\loai{Tính thời gian hoạt động}
\begin{vidu}%[3P7-5.2G][Loai1]
Một tàu ngầm hạt nhân có công suất 160 kW, dùng năng lượng phân hạch của hạt nhân U235 với hiệu suất $20\%$.  Lấy $N_A=6,{02.10}^{23}\ mol^{- 1}$ và khối lượng mol của \Hn{92}{235}{U} là $235\ g/mol$. Trung bình mỗi hạt nhân U235 phân hạch tỏa ra năng lượng 200\ MeV. Thời gian để tàu tiêu thụ hết 0,5 kg U235 nguyên chất \textbf{gần nhất} với giá trị nào sau đây?	
\choice
{529 ngày}
{756 ngày}
{\True 593 ngày}
{629 ngày}	
\loigiai{
\begin{itemize}
\item Năng lượng hạt nhân cần cung cấp: $Q=\dfrac{P_ct}{H}$.
\item Năng lượng của m gam hạt nhân phân hạch hết: $Q'=NW=\dfrac{m}{M}.N_A.W$.
\item Thời gian hoạt động của tàu: $Q=Q' \Leftrightarrow \dfrac{P_ct}{H}=\dfrac{mN_AW}{M} \Rightarrow t=\dfrac{HmN_AW}{MP_c}=51259574\ s \approx 592\ \text{ngày}$.
\end{itemize}
}
\end{vidu}
